\chapter{几何作图}\label{ch:g5:geometry}

用直尺、三角板、圆规这三件工具，再配合一点方法，
任何用文字描述的图形都能精确作出。本章
讲解作垂线、作\omterm{def:g4:geometry:def}{平行}线、作三角形以及
特殊\omterm{def:g4:geometry:polygon}{四边形}——先动手，后讲理论
（见\cref{ch:g6:lines,ch:g6:shapes}）。

\section{过一点作垂线与平行线}

\begin{method}[两种基本作图]\label{met:g5:geometry:basic}
\emph{作 $d$ 的垂线，过 $A$}：
沿 $d$ 滑动三角板，直到它的另一条直角边经过 $A$；画线；标出
\omterm{def:g3:shapes:rightangle}{直角}。

\emph{作 $d$ 的\omterm{def:g4:geometry:def}{平行}线，过 $A$}：把直尺紧靠三角板，
将直尺固定在 $d$ 上作导轨，把三角板沿直尺滑动至 $A$；
画线。（检验：两条直线之间的距离
处处相等。）
\end{method}

\begin{example}[连续作图]\label{ex:g5:geometry:chain}
许多图形都是这两种操作的连续组合。作长方形 $ABCD$，
其中 $AB = 5$ cm，$BC = 3$ cm：画 $[AB]$；在 $A$ 处作垂线；
在 $B$ 处作垂线；在同侧定出 $D$ 与 $C$，各长 $3$ cm；
连线。每一步都是一次基本作图。
\end{example}

\section{用圆规作三角形}

\begin{method}[已知三边作三角形]\label{met:g5:geometry:triangle}
作三边长分别为 $6$ cm、$5$ cm、$4$ cm 的三角形：
\begin{enumerate}
  \item 画最长的边 $[AB]$，其中 $AB = 6$ cm；
  \item 以 $A$ 为圆心、$5$ cm 为\omterm{def:g3:shapes:shapes}{半径}画弧；
  \item 以 $B$ 为圆心、$4$ cm 为\omterm{def:g3:shapes:shapes}{半径}画弧；
  \item 两弧的交点即第三个\omterm{def:g4:geometry:polygon}{顶点} $C$，连线。
\end{enumerate}
圆规保证了长度：第一段弧上每一点
$5$ cm 恰是它到 $A$ 的距离。
\end{method}

\begin{omfigure}
\begin{tikzpicture}[scale=0.95]
  \coordinate (A) at (0,0);
  \coordinate (B) at (6*0.72,0);
  \coordinate (C) at (2.7,2.376);
  \draw[thick] (A) -- (B);
  \draw[thick, omDef] (A) -- (C);
  \draw[thick, omThm] (B) -- (C);
  \draw[omDef, thin] (C) arc[start angle=41.3, delta angle=20,
    radius=3.6];
  \draw[omDef, thin] (C) arc[start angle=41.3, delta angle=-20,
    radius=3.6];
  \draw[omThm, thin] (C) arc[start angle=124.4, delta angle=20,
    radius=2.88];
  \draw[omThm, thin] (C) arc[start angle=124.4, delta angle=-20,
    radius=2.88];
  \fill (A) circle (1.5pt) node[below left, font=\small] {$A$};
  \fill (B) circle (1.5pt) node[below right, font=\small] {$B$};
  \fill (C) circle (1.5pt) node[above, font=\small] {$C$};
  \node[below, font=\small] at (2.16,0) {$6$ cm};
  \node[omDef, font=\small] at (0.85,1.45) {$5$ cm};
  \node[omThm, font=\small] at (3.9,1.35) {$4$ cm};
\end{tikzpicture}

{\small 两弧一交：用圆规作出三边为 $6$--$5$--$4$ 的三角形。
（若两条短边之和小于底边，两弧便不会相交——
这一点详见
\cref{ch:g7:triangles}。）}
\end{omfigure}

\begin{example}[轻松作出等腰与等边三角形]\label{ex:g5:geometry:special}
两弧所用圆规开口相同 $\Rightarrow$ 等腰三角形。
圆规开口等于底边 $\Rightarrow$ 等边三角形。圆规
并不量度，而是\emph{复制}长度，这样更好。
\end{example}

\section{用作图认识特殊四边形}

\begin{method}[依次作出各类四边形]\label{met:g5:geometry:quadrilaterals}
\begin{enumerate}
  \item \emph{长方形}：两对\omterm{def:g4:geometry:def}{平行}线相交成\omterm{def:g3:shapes:rightangle}{直角}
——或仿照\cref{ex:g5:geometry:chain}连续作图；
  \item \emph{正方形}：四条边长度相同的长方形
——用圆规搬运同一长度；
  \item \emph{菱形}：四条边都相等——先画一条边，再
        以它的两端为圆心、用相同开口画两段弧，定出
        另外两个\omterm{def:g4:geometry:polygon}{顶点}（菱形是两个等腰三角形
        沿一条对角线拼成的）。
\end{enumerate}
作图之后，务必用工具\emph{检验}：三角板
靠在各角上，圆规量各边。
\end{method}

\begin{example}[由描述辨认图形]\label{ex:g5:geometry:recognize}
“一个\omterm{def:g4:geometry:polygon}{四边形}有四条 $4$ cm 的边和一个\omterm{def:g3:shapes:rightangle}{直角}”——
试着作出它：四条相等的边决定了它是菱形，而将
一个角压成\omterm{def:g3:shapes:rightangle}{直角}，其余各角随之都成\omterm{def:g3:shapes:rightangle}{直角}：结果是
正方形。有些描述看似平常，实则描述了更特殊的
图形。（为什么一个\omterm{def:g3:shapes:rightangle}{直角}会使四个角都成\omterm{def:g3:shapes:rightangle}{直角}，
将在\cref{ch:g7:central}中说明。）
\end{example}

\section{作图中的圆}

\begin{example}[把圆当作工具]\label{ex:g5:geometry:circle}
“求所有距离为 $3$ cm 的点，距离自 $A$ 量起”：即以 $A$ 为圆心、
\omterm{def:g3:shapes:shapes}{半径}为 $3$ cm 的圆。“求所有距离为 $3$ cm 的点，距离自 $A$ 量起，\emph{且}
距离为 $4$ cm、自 $B$ 量起的点”：画出两个圆，它们的交点（两个、一个或
零个）即所求。寻宝图与三角形作图
用的正是这一思想。
\end{example}

\section{练习}

\begin{exercise}[$\star$]\label{exo:g5:geometry:1}
画一条直线 $d$，取一点 $A$，使其到直线的距离约为 $3$ cm。
作与 $d$ \omterm{def:g4:geometry:def}{垂直}且过 $A$ 的直线，再作与 $d$ \omterm{def:g4:geometry:def}{平行}的直线，
过 $A$ 的直线，并标出\omterm{def:g3:shapes:rightangle}{直角}。
\end{exercise}

\begin{exercise}[$\star$]\label{exo:g5:geometry:2}
作长方形 $ABCD$，使 $AB = 7$ cm，$BC = 4$ cm，
按\cref{ex:g5:geometry:chain}的方法作图，用三角尺检验
第四个角：若作图准确，该角
自然为\omterm{def:g3:shapes:rightangle}{直角}。
\end{exercise}

\begin{exercise}[$\star$]\label{exo:g5:geometry:3}
作三角形，三边长分别为 $7$ cm、$6$ cm 和 $3$ cm，
（见\cref{met:g5:geometry:triangle}）。
\end{exercise}

\begin{exercise}[$\star$]\label{exo:g5:geometry:4}
作等腰三角形，底边长 $5$ cm，腰长
$7$ cm；再作边长为 $6$ cm 的等边三角形。
\end{exercise}

\begin{exercise}[$\star$]\label{exo:g5:geometry:5}
作边长为 $5$ cm 的正方形。哪些工具能保证
\omterm{def:g3:shapes:rightangle}{直角}？哪些能保证边相等？
\end{exercise}

\begin{exercise}[$\star$]\label{exo:g5:geometry:6}
作边长为 $4$ cm、且不是正方形的菱形（第一个
角的张口可自选）。用圆规检验四条边
是否相等。
\end{exercise}

\begin{exercise}[$\star$]\label{exo:g5:geometry:7}
画两点 $A$、$B$，使 $AB = 5$ cm。作出所有
距离为 $4$ cm（自 $A$ 量起）、$3$ cm（自 $B$ 量起）的点，
共有几个？
\end{exercise}

\begin{exercise}[$\star$]\label{exo:g5:geometry:8}
写出一份作图程序（按序号列出步骤，并为各点命名），用于作
$6$ cm $\times$ $2$ cm 的长方形，上方接一个等腰三角形，
两腰长 $4$ cm（形如一支铅笔）。请一位同学
依此程序作图。
\end{exercise}

\begin{exercise}[$\star\star$]\label{exo:g5:geometry:9}
尝试作三边长分别为 $8$ cm、$3$ cm 和 $4$ cm 的三角形。
圆弧会出现什么情况？用一句话说明这个
三角形为何不存在。
\end{exercise}

\begin{exercise}[$\star\star$]\label{exo:g5:geometry:10}
作菱形，使其两条对角线长分别为 $6$ cm 和 $4$ cm，
已知菱形的对角线在中点处互相成\omterm{def:g3:shapes:rightangle}{直角}相交：
应\emph{先}画出对角线，再
连接它们的四个端点。
\end{exercise}

\begin{exercise}[$\star\star$]\label{exo:g5:geometry:11}
一只山羊拴在木桩 $A$ 上，绳长 $4$ m，
平地上有一条直线形栅栏，离 $3$ m 远（距离自 $A$ 量起）。画出平面图
（比例尺为每 1 cm 代表 1 m）：山羊能到达的区域是哪部分？
栅栏在何处限制了它？
\end{exercise}
